\section{Inconsistency from aggregate labels}
\label{sec:imp-mv}

%Our theoretical analysis begins with a few simple observations that at least
%suggest the importance of using non-aggregated labels: calibration
%for labeler uncertainty is, in a
%sense, impossible with only aggregated data, which is in strong contrast to
%algorithms that use non-aggregated data.
%
To make clear the impossibility of estimation without some knowledge of the
link and number $m$ of labelers, we show that there exist distributions with
different linear classifiers $\theta\opt$ and $\wb{\theta}$ generating the
same data distribution.  Consider a generalized linear model for binary
classification, with link  $\sigma\opt \in \fsymlink$, $\theta\opt
\in \R^d$, define the model \eqref{eq:general-calibration-model} with
$\Prb_{\sigma\opt, \theta\opt}(Y = y \mid X = x) = \sigma\opt ( y \<
\theta\opt, x\>)$, and let $X_1, X_2, \ldots, X_n \simiid \normal(0, I_d)$,
where labels follow $Y_{ij}\simiid
\Prb_{\sigma\opt,\theta\opt}(\cdot \mid X_i)$, $j = 1, \ldots, m$;
$\majY_i = \maj(Y_{i1}, \ldots, Y_{im})$ denotes the majority-vote.
Then
there always exists a calibration function $\wb{\sigma}$ and label size
$\wb{m}$ such that $\majY_i$ has identical distribution under both the
model~\eqref{eq:general-calibration-model} and with $\wb{\sigma},
\wb{m}$ replacing $\sigma\opt$ and $m$.  Letting
$\P_{(X,\majY)}^{\sigma\opt, \theta\opt, m}$ denote the induced distribution
on $(X_i, \majY_i)$, we have the following formal statement, which we prove
in Appendix~\ref{proof:impossibility}.

\begin{proposition} \label{prop:impossibility}
  Suppose $\sigma\opt \in \fsymlink$ satisfies $\sigma\opt(t) > 0$ for all $ t > 0$. For any $\wb{\theta} \in \R^d$ such that $\wb{\theta}/\ltwo{\wb{\theta}} = \theta\opt/\ltwo{\theta\opt}$ and any positive integer $\wb{m}$, there is another link function $\wb{\sigma}$ such that 
  \begin{align*}
    \P_{(X,\majY)}^{\sigma\opt, \theta\opt, m} \stackrel{\textup{dist}}{=}
    \P_{(X,\majY)}^{\wb{\sigma}, \wb{\theta}, \wb{m}}.
  \end{align*}
\end{proposition}

%\cc{A comment for this section is required. What are we doing here? We basically say if the non-aggregated data is completely not revealed to us, it is impossible to calibrate. This gives rise to the study of learning from non-aggregated data as it should carry more information, and the most important of which---whether it is possible to achieve calibrated estimation.  Although It is possible to achieve partial recovery only knowing the number of labelers $m$ with some additional assumptions, but we will focus on the regime learning from revealed annotations for all labelers. }
